% --- Archon physics formalization source begin ---
% archon:physics
% archon:covers PhyXMiniProblems/problem_phyx_mini_0317.lean
% archon:source-report reports/phyx_mini/problem_phyx_mini_0317.source.json

\chapter{Physics problem phyx\_mini\_0317}
\label{ch:PhyXMiniProblems_problem_phyx_mini_0317}

\paragraph{Problem source.}
\#\# Physical scenario

Two atmospheric sound sources \$A\$ and \$B\$ emit isotropically at constant power. The sound levels \$\textbackslash{}beta\$ of their emissions are plotted in figure versus the radial distance \$r\$ from the sources. The vertical axis scale is set by \$\textbackslash{}beta\_1 = 85.0\$ dB and \$\textbackslash{}beta\_2 = 65.0\$ dB.

\#\# Question

What is the sound level difference at \$r = 10\$ m?

\#\# Answer choices

- A: A: 3.5 dB

- B: B: 4.0 dB

- C: C: 4.5 dB

- D: D: 5.0 dB

\#\# Figure caption (auxiliary; use the image as primary evidence)

The image is a graph depicting two curves labeled 'A' and 'B'. The x-axis represents distance \textbackslash{}( r \textbackslash{}) in meters, ranging from 100 to 1000. The y-axis represents \textbackslash{}( \textbackslash{}beta \textbackslash{}) in decibels (dB), ranging from \textbackslash{}( \textbackslash{}beta\_2 \textbackslash{}) to \textbackslash{}( \textbackslash{}beta\_1 \textbackslash{}), with \textbackslash{}( \textbackslash{}beta\_1 \textbackslash{}) being the higher value. Both curves show a decreasing trend from left to right, indicating a reduction in \textbackslash{}( \textbackslash{}beta \textbackslash{}) as \textbackslash{}( r \textbackslash{}) increases. Curve 'A' starts at \textbackslash{}( \textbackslash{}beta\_1 \textbackslash{}) and decreases more steeply compared to curve 'B', which starts slightly lower than curve 'A'. The graph is on a grid background, providing a clear visual of the trend and relationship between the variables.

\#\# Dataset metadata

Sound; Physical Model Grounding Reasoning, Spatial Relation Reasoning

\paragraph{Recorded answer/context.}
D: D: 5.0 dB

\paragraph{Figure/image path.}
/root/proposal\_for\_physic/science-mango/phyx\_mini\_run/phyx\_data/test\_image/317.png

\paragraph{Formalization target.}
create a compiling Lean file with sorry bodies at `PhyXMiniProblems/problem\_phyx\_mini\_0317.lean`. The Lean declarations must preserve the physical quantities, dimensions or dimensional roles, figure labels, governing-law hypotheses, and final relation expressed by this problem.
Use Mathlib/Physlib names found through LeanExplore where available. If a physics API is missing, introduce faithful local abstractions rather than scalar placeholder aliases.

\begin{theorem}[Physics formalization target]
\label{thm:physics:phyx_mini_0317:target}
\lean{PhyXMiniProblems.ProblemPhyXMini0317.problem_phyx_mini_0317}
\uses{def:physics:phyx-mini-0317:phyxminiproblems-problemphyxmini0317-atmospherictwosourcesetup, def:physics:phyx-mini-0317:phyxminiproblems-problemphyxmini0317-soundleveldifferenceindecibels, def:physics:phyx-mini-0317:phyxminiproblems-problemphyxmini0317-matchesatmosphericsourcescenario, def:physics:phyx-mini-0317:phyxminiproblems-problemphyxmini0317-matchesproblemdistancedata, def:physics:phyx-mini-0317:phyxminiproblems-problemphyxmini0317-matchessuppliedsoundlevelgraph, def:physics:phyx-mini-0317:phyxminiproblems-problemphyxmini0317-usesstandardsoundintensityreference, def:physics:phyx-mini-0317:phyxminiproblems-problemphyxmini0317-hasphysicalacousticparameters, def:physics:phyx-mini-0317:phyxminiproblems-problemphyxmini0317-satisfiesisotropicinversesquarelaw, def:physics:phyx-mini-0317:phyxminiproblems-problemphyxmini0317-recordedanswerchoice, def:physics:phyx-mini-0317:phyxminiproblems-problemphyxmini0317-isclosestdisplayedanswer}
The assigned autoformalize agent should translate this physics problem into Lean declarations in the covered file, with theorem and lemma proof bodies written as `by sorry`.
\end{theorem}
\begin{proof}
This is an autoformalization task, not a proof task. Produce statements that can later be proved by the physics prover without weakening the physical model.
\end{proof}

% --- Archon declaration topology begin ---
\section{Declaration topology}
\begin{definition}[Length Magnitude]
  \label{def:physics:phyx-mini-0317:phyxminiproblems-problemphyxmini0317-lengthmagnitude}
  \lean{PhyXMiniProblems.ProblemPhyXMini0317.LengthMagnitude}
  A nonnegative physical radial distance from a sound source
\end{definition}
\begin{proof}
  This is the declared model definition of Length Magnitude.
\end{proof}
\begin{definition}[Acoustic Power Quantity]
  \label{def:physics:phyx-mini-0317:phyxminiproblems-problemphyxmini0317-acousticpowerquantity}
  \lean{PhyXMiniProblems.ProblemPhyXMini0317.AcousticPowerQuantity}
  Acoustic power has SI unit watt and dimension mass * length\textasciicircum{}2 / time\textasciicircum{}3
\end{definition}
\begin{proof}
  This is the declared model definition of Acoustic Power Quantity.
\end{proof}
\begin{definition}[Acoustic Intensity Quantity]
  \label{def:physics:phyx-mini-0317:phyxminiproblems-problemphyxmini0317-acousticintensityquantity}
  \lean{PhyXMiniProblems.ProblemPhyXMini0317.AcousticIntensityQuantity}
  Acoustic intensity is power per area, with SI unit watt per square metre and dimension mass / time\textasciicircum{}3
\end{definition}
\begin{proof}
  This is the declared model definition of Acoustic Intensity Quantity.
\end{proof}
\begin{definition}[distance Readout]
  \label{def:physics:phyx-mini-0317:phyxminiproblems-problemphyxmini0317-distancereadout}
  \lean{PhyXMiniProblems.ProblemPhyXMini0317.distanceReadout}
  \uses{def:physics:phyx-mini-0317:phyxminiproblems-problemphyxmini0317-lengthmagnitude}
  Read a physical distance in a coherent choice of units
\end{definition}
\begin{proof}
  This is the declared model definition of distance Readout.
\end{proof}
\begin{definition}[acoustic Power Readout]
  \label{def:physics:phyx-mini-0317:phyxminiproblems-problemphyxmini0317-acousticpowerreadout}
  \lean{PhyXMiniProblems.ProblemPhyXMini0317.acousticPowerReadout}
  \uses{def:physics:phyx-mini-0317:phyxminiproblems-problemphyxmini0317-acousticpowerquantity}
  Read an acoustic power in a coherent choice of units
\end{definition}
\begin{proof}
  This is the declared model definition of acoustic Power Readout.
\end{proof}
\begin{definition}[acoustic Intensity Readout]
  \label{def:physics:phyx-mini-0317:phyxminiproblems-problemphyxmini0317-acousticintensityreadout}
  \lean{PhyXMiniProblems.ProblemPhyXMini0317.acousticIntensityReadout}
  \uses{def:physics:phyx-mini-0317:phyxminiproblems-problemphyxmini0317-acousticintensityquantity}
  Read an acoustic intensity in a coherent choice of units
\end{definition}
\begin{proof}
  This is the declared model definition of acoustic Intensity Readout.
\end{proof}
\begin{definition}[distance In Meters]
  \label{def:physics:phyx-mini-0317:phyxminiproblems-problemphyxmini0317-distanceinmeters}
  \lean{PhyXMiniProblems.ProblemPhyXMini0317.distanceInMeters}
  \uses{def:physics:phyx-mini-0317:phyxminiproblems-problemphyxmini0317-lengthmagnitude, def:physics:phyx-mini-0317:phyxminiproblems-problemphyxmini0317-distancereadout}
  SI metre readout of a radial distance
\end{definition}
\begin{proof}
  This is the declared model definition of distance In Meters.
\end{proof}
\begin{definition}[acoustic Power In Watts]
  \label{def:physics:phyx-mini-0317:phyxminiproblems-problemphyxmini0317-acousticpowerinwatts}
  \lean{PhyXMiniProblems.ProblemPhyXMini0317.acousticPowerInWatts}
  \uses{def:physics:phyx-mini-0317:phyxminiproblems-problemphyxmini0317-acousticpowerquantity, def:physics:phyx-mini-0317:phyxminiproblems-problemphyxmini0317-acousticpowerreadout}
  SI watt readout of an emitted acoustic power
\end{definition}
\begin{proof}
  This is the declared model definition of acoustic Power In Watts.
\end{proof}
\begin{definition}[acoustic Intensity In Watts Per Square Meter]
  \label{def:physics:phyx-mini-0317:phyxminiproblems-problemphyxmini0317-acousticintensityinwattspersquaremeter}
  \lean{PhyXMiniProblems.ProblemPhyXMini0317.acousticIntensityInWattsPerSquareMeter}
  \uses{def:physics:phyx-mini-0317:phyxminiproblems-problemphyxmini0317-acousticintensityquantity, def:physics:phyx-mini-0317:phyxminiproblems-problemphyxmini0317-acousticintensityreadout}
  SI watt-per-square-metre readout of an acoustic intensity
\end{definition}
\begin{proof}
  This is the declared model definition of acoustic Intensity In Watts Per Square Meter.
\end{proof}
\begin{definition}[Source Label]
  \label{def:physics:phyx-mini-0317:phyxminiproblems-problemphyxmini0317-sourcelabel}
  \lean{PhyXMiniProblems.ProblemPhyXMini0317.SourceLabel}
  The two curves and atmospheric sound sources labeled in the figure
\end{definition}
\begin{proof}
  This is the declared model definition of Source Label.
\end{proof}
\begin{definition}[Radial Position]
  \label{def:physics:phyx-mini-0317:phyxminiproblems-problemphyxmini0317-radialposition}
  \lean{PhyXMiniProblems.ProblemPhyXMini0317.RadialPosition}
  Named radial positions. The requested10m point comes from the question; the remaining positions are the labeled marks on the supplied graph
\end{definition}
\begin{proof}
  This is the declared model definition of Radial Position.
\end{proof}
\begin{definition}[Acoustic Source Kind]
  \label{def:physics:phyx-mini-0317:phyxminiproblems-problemphyxmini0317-acousticsourcekind}
  \lean{PhyXMiniProblems.ProblemPhyXMini0317.AcousticSourceKind}
  Physical type of an idealized atmospheric acoustic source
\end{definition}
\begin{proof}
  This is the declared model definition of Acoustic Source Kind.
\end{proof}
\begin{definition}[Emission Regime]
  \label{def:physics:phyx-mini-0317:phyxminiproblems-problemphyxmini0317-emissionregime}
  \lean{PhyXMiniProblems.ProblemPhyXMini0317.EmissionRegime}
  Time dependence of a source's emitted acoustic power
\end{definition}
\begin{proof}
  This is the declared model definition of Emission Regime.
\end{proof}
\begin{definition}[Acoustic Medium]
  \label{def:physics:phyx-mini-0317:phyxminiproblems-problemphyxmini0317-acousticmedium}
  \lean{PhyXMiniProblems.ProblemPhyXMini0317.AcousticMedium}
  Propagation medium named in the physical scenario
\end{definition}
\begin{proof}
  This is the declared model definition of Acoustic Medium.
\end{proof}
\begin{definition}[Graph Axis Quantity]
  \label{def:physics:phyx-mini-0317:phyxminiproblems-problemphyxmini0317-graphaxisquantity}
  \lean{PhyXMiniProblems.ProblemPhyXMini0317.GraphAxisQuantity}
  \uses{def:physics:phyx-mini-0317:phyxminiproblems-problemphyxmini0317-soundlevelindecibels}
  Physical quantity represented by a graph axis
\end{definition}
\begin{proof}
  This is the declared model definition of Graph Axis Quantity.
\end{proof}
\begin{definition}[Sound Level Graph]
  \label{def:physics:phyx-mini-0317:phyxminiproblems-problemphyxmini0317-soundlevelgraph}
  \lean{PhyXMiniProblems.ProblemPhyXMini0317.SoundLevelGraph}
  \uses{def:physics:phyx-mini-0317:phyxminiproblems-problemphyxmini0317-sourcelabel, def:physics:phyx-mini-0317:phyxminiproblems-problemphyxmini0317-graphaxisquantity}
  Metadata and scalar readouts attached to the supplied graph. The two β fields are sound-level readouts in decibels, not physical intensities
\end{definition}
\begin{proof}
  This is the declared model definition of Sound Level Graph.
\end{proof}
\begin{definition}[vertical Grid Step In Decibels]
  \label{def:physics:phyx-mini-0317:phyxminiproblems-problemphyxmini0317-verticalgridstepindecibels}
  \lean{PhyXMiniProblems.ProblemPhyXMini0317.verticalGridStepInDecibels}
  \uses{def:physics:phyx-mini-0317:phyxminiproblems-problemphyxmini0317-soundlevelgraph}
  The size of one vertical grid step. In the image, β₁ and β₂ are separated by four equal grid intervals. This is a generic graph-scale conversion and does not mention the requested 10 m sound-level difference
\end{definition}
\begin{proof}
  This is the declared model definition of vertical Grid Step In Decibels.
\end{proof}
\begin{definition}[Atmospheric Two Source Setup]
  \label{def:physics:phyx-mini-0317:phyxminiproblems-problemphyxmini0317-atmospherictwosourcesetup}
  \lean{PhyXMiniProblems.ProblemPhyXMini0317.AtmosphericTwoSourceSetup}
  \uses{def:physics:phyx-mini-0317:phyxminiproblems-problemphyxmini0317-lengthmagnitude, def:physics:phyx-mini-0317:phyxminiproblems-problemphyxmini0317-acousticpowerquantity, def:physics:phyx-mini-0317:phyxminiproblems-problemphyxmini0317-acousticintensityquantity, def:physics:phyx-mini-0317:phyxminiproblems-problemphyxmini0317-sourcelabel, def:physics:phyx-mini-0317:phyxminiproblems-problemphyxmini0317-radialposition, def:physics:phyx-mini-0317:phyxminiproblems-problemphyxmini0317-acousticsourcekind, def:physics:phyx-mini-0317:phyxminiproblems-problemphyxmini0317-emissionregime, def:physics:phyx-mini-0317:phyxminiproblems-problemphyxmini0317-acousticmedium, def:physics:phyx-mini-0317:phyxminiproblems-problemphyxmini0317-soundlevelgraph}
  Physical setup for the two-source comparison. Intensities at all named radii remain independent fields until related by the governing inverse-square law. In particular, no intensity ratio or requested level difference is stored here
\end{definition}
\begin{proof}
  This is the declared model definition of Atmospheric Two Source Setup.
\end{proof}
\begin{definition}[sound Level In Decibels]
  \label{def:physics:phyx-mini-0317:phyxminiproblems-problemphyxmini0317-soundlevelindecibels}
  \lean{PhyXMiniProblems.ProblemPhyXMini0317.soundLevelInDecibels}
  \uses{def:physics:phyx-mini-0317:phyxminiproblems-problemphyxmini0317-acousticintensityinwattspersquaremeter, def:physics:phyx-mini-0317:phyxminiproblems-problemphyxmini0317-sourcelabel, def:physics:phyx-mini-0317:phyxminiproblems-problemphyxmini0317-radialposition, def:physics:phyx-mini-0317:phyxminiproblems-problemphyxmini0317-atmospherictwosourcesetup}
  Sound intensity level in decibels, using the standard general relation β = 10 log10 (I / I\_ref) and same-unit SI intensity readouts
\end{definition}
\begin{proof}
  This is the declared model definition of sound Level In Decibels.
\end{proof}
\begin{definition}[sound Level Difference In Decibels]
  \label{def:physics:phyx-mini-0317:phyxminiproblems-problemphyxmini0317-soundleveldifferenceindecibels}
  \lean{PhyXMiniProblems.ProblemPhyXMini0317.soundLevelDifferenceInDecibels}
  \uses{def:physics:phyx-mini-0317:phyxminiproblems-problemphyxmini0317-radialposition, def:physics:phyx-mini-0317:phyxminiproblems-problemphyxmini0317-atmospherictwosourcesetup, def:physics:phyx-mini-0317:phyxminiproblems-problemphyxmini0317-soundlevelindecibels}
  The positive magnitude of the level difference between curves A and B at a common radius. This definition contains no expected numerical answer
\end{definition}
\begin{proof}
  This is the declared model definition of sound Level Difference In Decibels.
\end{proof}
\begin{definition}[Matches Atmospheric Source Scenario]
  \label{def:physics:phyx-mini-0317:phyxminiproblems-problemphyxmini0317-matchesatmosphericsourcescenario}
  \lean{PhyXMiniProblems.ProblemPhyXMini0317.MatchesAtmosphericSourceScenario}
  \uses{def:physics:phyx-mini-0317:phyxminiproblems-problemphyxmini0317-sourcelabel, def:physics:phyx-mini-0317:phyxminiproblems-problemphyxmini0317-atmospherictwosourcesetup}
  Both labeled atmospheric sources are isotropic point sources emitting at constant (not necessarily equal) power through the atmosphere
\end{definition}
\begin{proof}
  This is the declared model definition of Matches Atmospheric Source Scenario.
\end{proof}
\begin{definition}[Matches Problem Distance Data]
  \label{def:physics:phyx-mini-0317:phyxminiproblems-problemphyxmini0317-matchesproblemdistancedata}
  \lean{PhyXMiniProblems.ProblemPhyXMini0317.MatchesProblemDistanceData}
  \uses{def:physics:phyx-mini-0317:phyxminiproblems-problemphyxmini0317-distanceinmeters, def:physics:phyx-mini-0317:phyxminiproblems-problemphyxmini0317-atmospherictwosourcesetup}
  The numerical distance readouts in the question and on the horizontal graph axis. Recording 10 m identifies the requested position but says nothing about either source's level there
\end{definition}
\begin{proof}
  This is the declared model definition of Matches Problem Distance Data.
\end{proof}
\begin{definition}[Matches Supplied Sound Level Graph]
  \label{def:physics:phyx-mini-0317:phyxminiproblems-problemphyxmini0317-matchessuppliedsoundlevelgraph}
  \lean{PhyXMiniProblems.ProblemPhyXMini0317.MatchesSuppliedSoundLevelGraph}
  \uses{def:physics:phyx-mini-0317:phyxminiproblems-problemphyxmini0317-verticalgridstepindecibels, def:physics:phyx-mini-0317:phyxminiproblems-problemphyxmini0317-atmospherictwosourcesetup, def:physics:phyx-mini-0317:phyxminiproblems-problemphyxmini0317-soundlevelindecibels}
  Primary-image evidence from the plotted graph: the horizontal axis is radial distance in metres and the vertical axis is sound level in decibels; curves A and B are both visible; β₁ = 85 dB and β₂ = 65 dB; at 100 m, curve B meets β₁, while curve A is one of the four equal vertical grid steps higher; and both plotted curves decrease across the labeled graph range. These are graph readouts at 100 m, 500 m, and 1000 m; no level or level difference at the requested 10 m position occurs in this structure
\end{definition}
\begin{proof}
  This is the declared model definition of Matches Supplied Sound Level Graph.
\end{proof}
\begin{definition}[Uses Standard Sound Intensity Reference]
  \label{def:physics:phyx-mini-0317:phyxminiproblems-problemphyxmini0317-usesstandardsoundintensityreference}
  \lean{PhyXMiniProblems.ProblemPhyXMini0317.UsesStandardSoundIntensityReference}
  \uses{def:physics:phyx-mini-0317:phyxminiproblems-problemphyxmini0317-acousticintensityinwattspersquaremeter, def:physics:phyx-mini-0317:phyxminiproblems-problemphyxmini0317-atmospherictwosourcesetup}
  The conventional common reference intensity for sound intensity level is 10⁻¹² W/m². This calibration is independent of the source and radius and does not determine the requested difference, which cancels the common reference
\end{definition}
\begin{proof}
  This is the declared model definition of Uses Standard Sound Intensity Reference.
\end{proof}
\begin{definition}[Has Physical Acoustic Parameters]
  \label{def:physics:phyx-mini-0317:phyxminiproblems-problemphyxmini0317-hasphysicalacousticparameters}
  \lean{PhyXMiniProblems.ProblemPhyXMini0317.HasPhysicalAcousticParameters}
  \uses{def:physics:phyx-mini-0317:phyxminiproblems-problemphyxmini0317-distanceinmeters, def:physics:phyx-mini-0317:phyxminiproblems-problemphyxmini0317-acousticpowerinwatts, def:physics:phyx-mini-0317:phyxminiproblems-problemphyxmini0317-acousticintensityinwattspersquaremeter, def:physics:phyx-mini-0317:phyxminiproblems-problemphyxmini0317-sourcelabel, def:physics:phyx-mini-0317:phyxminiproblems-problemphyxmini0317-radialposition, def:physics:phyx-mini-0317:phyxminiproblems-problemphyxmini0317-atmospherictwosourcesetup}
  Strict positivity and nondegeneracy of all physical acoustic quantities
\end{definition}
\begin{proof}
  This is the declared model definition of Has Physical Acoustic Parameters.
\end{proof}
\begin{definition}[Satisfies Isotropic Inverse Square Law]
  \label{def:physics:phyx-mini-0317:phyxminiproblems-problemphyxmini0317-satisfiesisotropicinversesquarelaw}
  \lean{PhyXMiniProblems.ProblemPhyXMini0317.SatisfiesIsotropicInverseSquareLaw}
  \uses{def:physics:phyx-mini-0317:phyxminiproblems-problemphyxmini0317-distancereadout, def:physics:phyx-mini-0317:phyxminiproblems-problemphyxmini0317-acousticpowerreadout, def:physics:phyx-mini-0317:phyxminiproblems-problemphyxmini0317-acousticintensityreadout, def:physics:phyx-mini-0317:phyxminiproblems-problemphyxmini0317-sourcelabel, def:physics:phyx-mini-0317:phyxminiproblems-problemphyxmini0317-radialposition, def:physics:phyx-mini-0317:phyxminiproblems-problemphyxmini0317-atmospherictwosourcesetup}
  The spherical inverse-square intensity law for each isotropic source. For every named radius and coherent unit system, the emitted power equals 4 * π * r\textasciicircum{}2 * I. This generic law contains neither a particular power ratio nor the requested 5 dB conclusion
\end{definition}
\begin{proof}
  This is the declared model definition of Satisfies Isotropic Inverse Square Law.
\end{proof}
\begin{lemma}[sound Level Difference independent Of Radius]
  \label{lem:physics:phyx-mini-0317:phyxminiproblems-problemphyxmini0317-soundleveldifference-independentofradius}
  \lean{PhyXMiniProblems.ProblemPhyXMini0317.soundLevelDifference_independentOfRadius}
  \uses{def:physics:phyx-mini-0317:phyxminiproblems-problemphyxmini0317-radialposition, def:physics:phyx-mini-0317:phyxminiproblems-problemphyxmini0317-atmospherictwosourcesetup, def:physics:phyx-mini-0317:phyxminiproblems-problemphyxmini0317-soundleveldifferenceindecibels, def:physics:phyx-mini-0317:phyxminiproblems-problemphyxmini0317-hasphysicalacousticparameters, def:physics:phyx-mini-0317:phyxminiproblems-problemphyxmini0317-satisfiesisotropicinversesquarelaw}
  Because both sources obey the inverse-square law, their intensity ratio and therefore their decibel-level difference are independent of the common radial position. This is derived from the governing law, not assumed in the setup
\end{lemma}
\begin{proof}
  The conclusion follows from the governing relations and figure readouts named in the statement of sound Level Difference independent Of Radius.
\end{proof}
\begin{lemma}[graph Sound Level Difference At100m]
  \label{lem:physics:phyx-mini-0317:phyxminiproblems-problemphyxmini0317-graphsoundleveldifferenceat100m}
  \lean{PhyXMiniProblems.ProblemPhyXMini0317.graphSoundLevelDifferenceAt100m}
  \uses{def:physics:phyx-mini-0317:phyxminiproblems-problemphyxmini0317-atmospherictwosourcesetup, def:physics:phyx-mini-0317:phyxminiproblems-problemphyxmini0317-soundleveldifferenceindecibels, def:physics:phyx-mini-0317:phyxminiproblems-problemphyxmini0317-matchessuppliedsoundlevelgraph}
  The graph scale and the two curve readings at 100 m give a one-grid-step separation of (85 - 65) / 4 = 5 dB. This remains a derived statement
\end{lemma}
\begin{proof}
  The conclusion follows from the governing relations and figure readouts named in the statement of graph Sound Level Difference At100m.
\end{proof}
\begin{definition}[Answer Choice]
  \label{def:physics:phyx-mini-0317:phyxminiproblems-problemphyxmini0317-answerchoice}
  \lean{PhyXMiniProblems.ProblemPhyXMini0317.AnswerChoice}
  Labels of the four multiple-choice answers printed with the problem
\end{definition}
\begin{proof}
  This is the declared model definition of Answer Choice.
\end{proof}
\begin{definition}[decibels]
  \label{def:physics:phyx-mini-0317:phyxminiproblems-problemphyxmini0317-answerchoice-decibels}
  \lean{PhyXMiniProblems.ProblemPhyXMini0317.AnswerChoice.decibels}
  \uses{def:physics:phyx-mini-0317:phyxminiproblems-problemphyxmini0317-answerchoice}
  Sound-level difference in decibels printed beside each answer label
\end{definition}
\begin{proof}
  This is the declared model definition of decibels.
\end{proof}
\begin{definition}[recorded Answer Choice]
  \label{def:physics:phyx-mini-0317:phyxminiproblems-problemphyxmini0317-recordedanswerchoice}
  \lean{PhyXMiniProblems.ProblemPhyXMini0317.recordedAnswerChoice}
  \uses{def:physics:phyx-mini-0317:phyxminiproblems-problemphyxmini0317-answerchoice}
  The answer label recorded in the supplied dataset
\end{definition}
\begin{proof}
  This is the declared model definition of recorded Answer Choice.
\end{proof}
\begin{definition}[Is Closest Displayed Answer]
  \label{def:physics:phyx-mini-0317:phyxminiproblems-problemphyxmini0317-isclosestdisplayedanswer}
  \lean{PhyXMiniProblems.ProblemPhyXMini0317.IsClosestDisplayedAnswer}
  \uses{def:physics:phyx-mini-0317:phyxminiproblems-problemphyxmini0317-answerchoice}
  A displayed answer is uniquely closest to the physically derived sound-level difference. This generic comparison does not privilege choice D
\end{definition}
\begin{proof}
  This is the declared model definition of Is Closest Displayed Answer.
\end{proof}
% --- Archon declaration topology end ---
% --- Archon physics formalization source end ---
